%% 00-abstract.tex
%%
%% Length discipline: under 250 words. An earlier draft ran to 489 and read as a
%% table of contents. What was cut was the enumeration of everything the paper
%% covers; what was kept is the problem, the abstraction, the claim, and the
%% honest statement of what is not established. `make check' asserts the count.

\begin{abstract}
Production LLM-agent systems are well observed, while release governance remains
uneven across the resources they depend on. Tracing tells an
operator that a prompt regressed; the toolchain around it says little about what to
change, whether the change is better, or how to undo it. Prompt-lifecycle platforms
now close part of this for prompts --- versions, a named pointer resolved at call
time, rollback by re-pointing --- but not for the workflows, tool definitions, MCP
servers, and retrieval corpora an agent equally depends on, and not with the
evaluation result acting as a precondition rather than an annotation.

We present \sysbare{}, a layered control plane for the lifecycle of AI-agent
resources. Its organizing abstraction is the \emph{\gasset}: a typed record
carrying version history, an authority model, an evidence base, and --- where its
family supports one --- a \livetarget{} that client code resolves at call time.
\sys spans \statFamiliesW{} \families{} in \statPyLoc{} lines of Python and
\statUiLoc{} of TypeScript, reusing \mlflow{} rather than rebuilding it.

Layering makes shared capabilities available; explicit per-family wiring determines
the guarantees. We report that surface in full. Prompt authoring is non-live, and
every routed prompt-alias mutation uses a durable intent, idempotent provider operation,
and observable reconciliation state; equivalent crash handling is not yet uniform
across other families. The quantitative evaluation is specified but not executed;
every unpopulated result is marked rather than estimated.
\end{abstract}
